% --- Archon physics formalization source begin ---
% archon:physics
% archon:covers PhyXMiniProblems/problem_phyx_mini_0520.lean
% archon:source-report reports/phyx_mini/problem_phyx_mini_0520.source.json

\chapter{Physics problem phyx\_mini\_0520}
\label{ch:PhyXMiniProblems_problem_phyx_mini_0520}

\paragraph{Problem source.}
\#\# Physical scenario

In a simple model for a radioactive nucleus, an alpha particle \} \textbackslash{}( m = 6.64 \textbackslash{}times 10\textasciicircum{}\{-27\} \textbackslash{}text\{ kg\} \textbackslash{}) \textbackslash{}text\{ is trapped by a square barrier that has width \} \textbackslash{}( 2.0 \textbackslash{}text\{ fm\} \textbackslash{}) \textbackslash{}text\{ and height \} \textbackslash{}( 30.0 \textbackslash{}text\{ MeV\} \textbackslash{})

\#\# Question

What is the tunneling probability when the alpha particle encounters the barrier if its kinetic energy is \textbackslash{}( 1.0 \textbackslash{}text\{ MeV\} \textbackslash{}) below the top of the barrier as shown in figure?

\#\# Answer choices

- A: A: 0.065

- B: B: 0.042

- C: C: 0.084

- D: D: 0.090

\#\# Figure caption (auxiliary; use the image as primary evidence)

The image is a graph with the following elements:

1. **Axes**: 
   - The vertical axis is labeled \textbackslash{}( U(r) \textbackslash{}).
   - The horizontal axis is labeled \textbackslash{}( r \textbackslash{}).

2. **Potential Barrier**: 
   - A blue rectangular potential barrier is shown.
   - The barrier's height reaches up to the level marked by \textbackslash{}( U\_0 \textbackslash{}) on the vertical axis.
   - The width of the barrier is labeled as \textbackslash{}( 2.0 \textbackslash{}, \textbackslash{}text\{fm\} \textbackslash{}) at the base on the horizontal axis.

3. **Energy Level**:
   - A dashed brown line runs horizontally across the graph, labeled \textbackslash{}( E \textbackslash{}).
   - Above this, another dotted line labeled \textbackslash{}( 1.0 \textbackslash{}, \textbackslash{}text\{MeV\} \textbackslash{}) indicates the separation from \textbackslash{}( E \textbackslash{}) to \textbackslash{}( U\_0 \textbackslash{}).

4. **Labels and Annotations**:
   - The base of the potential barrier starts at the origin labeled \textbackslash{}( O \textbackslash{}).
   - The gap between the energy line \textbackslash{}( E \textbackslash{}) and \textbackslash{}( U\_0 \textbackslash{}) is annotated with an arrow and labeled \textbackslash{}( 1.0 \textbackslash{}, \textbackslash{}text\{MeV\} \textbackslash{}). 

This graph likely represents a quantum mechanics scenario, such as a particle encountering a potential barrier.

\#\# Dataset metadata

Quantum Phenomena; Physical Model Grounding Reasoning, Numerical Reasoning

\paragraph{Recorded answer/context.}
D: D: 0.090

\paragraph{Figure/image path.}
/root/proposal\_for\_physic/science-mango/phyx\_mini\_run/phyx\_data/test\_image/520.png

\paragraph{Formalization target.}
create a compiling Lean file with sorry bodies at `PhyXMiniProblems/problem\_phyx\_mini\_0520.lean`. The Lean declarations must preserve the physical quantities, dimensions or dimensional roles, figure labels, governing-law hypotheses, and final relation expressed by this problem.
Use Mathlib/Physlib names found through LeanExplore where available. If a physics API is missing, introduce faithful local abstractions rather than scalar placeholder aliases.

\begin{theorem}[Physics formalization target]
\label{thm:physics:phyx_mini_0520:target}
\lean{PhyXMiniProblems.ProblemPhyXMini0520.problem_phyx_mini_0520}
\uses{def:physics:phyx-mini-0520:phyxminiproblems-problemphyxmini0520-energyinmegaelectronvolts, def:physics:phyx-mini-0520:phyxminiproblems-problemphyxmini0520-alpharectangularbarriersetup, def:physics:phyx-mini-0520:phyxminiproblems-problemphyxmini0520-matchesalphanucleusscenario, def:physics:phyx-mini-0520:phyxminiproblems-problemphyxmini0520-matchessuppliedbarrierfigure, def:physics:phyx-mini-0520:phyxminiproblems-problemphyxmini0520-matchesproblemreadouts, def:physics:phyx-mini-0520:phyxminiproblems-problemphyxmini0520-usesstandardreducedplanckconstant, def:physics:phyx-mini-0520:phyxminiproblems-problemphyxmini0520-matchesphyslibrectangularbarrier, def:physics:phyx-mini-0520:phyxminiproblems-problemphyxmini0520-hasphysicaltunnelingparameters, def:physics:phyx-mini-0520:phyxminiproblems-problemphyxmini0520-barrieropacity, def:physics:phyx-mini-0520:phyxminiproblems-problemphyxmini0520-obeysfiniterectangularbarriertunneling, def:physics:phyx-mini-0520:phyxminiproblems-problemphyxmini0520-roundstothreedecimalplaces, def:physics:phyx-mini-0520:phyxminiproblems-problemphyxmini0520-isuniquematchingtunnelingchoice, def:physics:phyx-mini-0520:phyxminiproblems-problemphyxmini0520-nodisplayedtunnelingchoicematches}
The assigned autoformalize agent should translate this physics problem into Lean declarations in the covered file, with theorem and lemma proof bodies written as `by sorry`.
\end{theorem}
\begin{proof}
This is an autoformalization task, not a proof task. Produce statements that can later be proved by the physics prover without weakening the physical model.
\end{proof}

% --- Archon declaration topology begin ---
\section{Declaration topology}
\begin{definition}[Length Quantity]
  \label{def:physics:phyx-mini-0520:phyxminiproblems-problemphyxmini0520-lengthquantity}
  \lean{PhyXMiniProblems.ProblemPhyXMini0520.LengthQuantity}
  A nonnegative, unit-independent physical length
\end{definition}
\begin{proof}
  This is the declared model definition of Length Quantity.
\end{proof}
\begin{definition}[Mass Quantity]
  \label{def:physics:phyx-mini-0520:phyxminiproblems-problemphyxmini0520-massquantity}
  \lean{PhyXMiniProblems.ProblemPhyXMini0520.MassQuantity}
  A nonnegative, unit-independent physical mass
\end{definition}
\begin{proof}
  This is the declared model definition of Mass Quantity.
\end{proof}
\begin{definition}[Energy Quantity]
  \label{def:physics:phyx-mini-0520:phyxminiproblems-problemphyxmini0520-energyquantity}
  \lean{PhyXMiniProblems.ProblemPhyXMini0520.EnergyQuantity}
  A signed, unit-independent physical energy
\end{definition}
\begin{proof}
  This is the declared model definition of Energy Quantity.
\end{proof}
\begin{definition}[action Dimension]
  \label{def:physics:phyx-mini-0520:phyxminiproblems-problemphyxmini0520-actiondimension}
  \lean{PhyXMiniProblems.ProblemPhyXMini0520.actionDimension}
  The dimension of action, mass * length\textasciicircum{}2 / time
\end{definition}
\begin{proof}
  This is the declared model definition of action Dimension.
\end{proof}
\begin{definition}[Action Quantity]
  \label{def:physics:phyx-mini-0520:phyxminiproblems-problemphyxmini0520-actionquantity}
  \lean{PhyXMiniProblems.ProblemPhyXMini0520.ActionQuantity}
  \uses{def:physics:phyx-mini-0520:phyxminiproblems-problemphyxmini0520-actiondimension}
  A nonnegative physical action, used for the reduced Planck constant
\end{definition}
\begin{proof}
  This is the declared model definition of Action Quantity.
\end{proof}
\begin{definition}[Wave Number Quantity]
  \label{def:physics:phyx-mini-0520:phyxminiproblems-problemphyxmini0520-wavenumberquantity}
  \lean{PhyXMiniProblems.ProblemPhyXMini0520.WaveNumberQuantity}
  A nonnegative physical wave number, carrying inverse-length dimension
\end{definition}
\begin{proof}
  This is the declared model definition of Wave Number Quantity.
\end{proof}
\begin{definition}[length Readout]
  \label{def:physics:phyx-mini-0520:phyxminiproblems-problemphyxmini0520-lengthreadout}
  \lean{PhyXMiniProblems.ProblemPhyXMini0520.lengthReadout}
  \uses{def:physics:phyx-mini-0520:phyxminiproblems-problemphyxmini0520-lengthquantity}
  Read a physical length in a selected length unit
\end{definition}
\begin{proof}
  This is the declared model definition of length Readout.
\end{proof}
\begin{definition}[length In Meters]
  \label{def:physics:phyx-mini-0520:phyxminiproblems-problemphyxmini0520-lengthinmeters}
  \lean{PhyXMiniProblems.ProblemPhyXMini0520.lengthInMeters}
  \uses{def:physics:phyx-mini-0520:phyxminiproblems-problemphyxmini0520-lengthquantity, def:physics:phyx-mini-0520:phyxminiproblems-problemphyxmini0520-lengthreadout}
  Metre readout of a physical length
\end{definition}
\begin{proof}
  This is the declared model definition of length In Meters.
\end{proof}
\begin{definition}[length In Femtometers]
  \label{def:physics:phyx-mini-0520:phyxminiproblems-problemphyxmini0520-lengthinfemtometers}
  \lean{PhyXMiniProblems.ProblemPhyXMini0520.lengthInFemtometers}
  \uses{def:physics:phyx-mini-0520:phyxminiproblems-problemphyxmini0520-lengthquantity, def:physics:phyx-mini-0520:phyxminiproblems-problemphyxmini0520-lengthreadout}
  Femtometre readout of a physical length
\end{definition}
\begin{proof}
  This is the declared model definition of length In Femtometers.
\end{proof}
\begin{definition}[mass In Kilograms]
  \label{def:physics:phyx-mini-0520:phyxminiproblems-problemphyxmini0520-massinkilograms}
  \lean{PhyXMiniProblems.ProblemPhyXMini0520.massInKilograms}
  \uses{def:physics:phyx-mini-0520:phyxminiproblems-problemphyxmini0520-massquantity}
  Kilogram readout of a physical mass
\end{definition}
\begin{proof}
  This is the declared model definition of mass In Kilograms.
\end{proof}
\begin{definition}[energy In Joules]
  \label{def:physics:phyx-mini-0520:phyxminiproblems-problemphyxmini0520-energyinjoules}
  \lean{PhyXMiniProblems.ProblemPhyXMini0520.energyInJoules}
  \uses{def:physics:phyx-mini-0520:phyxminiproblems-problemphyxmini0520-energyquantity}
  Joule readout of a physical energy
\end{definition}
\begin{proof}
  This is the declared model definition of energy In Joules.
\end{proof}
\begin{definition}[energy In Electron Volts]
  \label{def:physics:phyx-mini-0520:phyxminiproblems-problemphyxmini0520-energyinelectronvolts}
  \lean{PhyXMiniProblems.ProblemPhyXMini0520.energyInElectronVolts}
  \uses{def:physics:phyx-mini-0520:phyxminiproblems-problemphyxmini0520-energyquantity}
  Electron-volt readout of a physical energy
\end{definition}
\begin{proof}
  This is the declared model definition of energy In Electron Volts.
\end{proof}
\begin{definition}[energy In Mega Electron Volts]
  \label{def:physics:phyx-mini-0520:phyxminiproblems-problemphyxmini0520-energyinmegaelectronvolts}
  \lean{PhyXMiniProblems.ProblemPhyXMini0520.energyInMegaElectronVolts}
  \uses{def:physics:phyx-mini-0520:phyxminiproblems-problemphyxmini0520-energyquantity, def:physics:phyx-mini-0520:phyxminiproblems-problemphyxmini0520-energyinelectronvolts}
  Mega-electron-volt readout of a physical energy
\end{definition}
\begin{proof}
  This is the declared model definition of energy In Mega Electron Volts.
\end{proof}
\begin{definition}[action In Joule Seconds]
  \label{def:physics:phyx-mini-0520:phyxminiproblems-problemphyxmini0520-actioninjouleseconds}
  \lean{PhyXMiniProblems.ProblemPhyXMini0520.actionInJouleSeconds}
  \uses{def:physics:phyx-mini-0520:phyxminiproblems-problemphyxmini0520-actionquantity}
  Joule-second readout of a physical action
\end{definition}
\begin{proof}
  This is the declared model definition of action In Joule Seconds.
\end{proof}
\begin{definition}[wave Number In Inverse Meters]
  \label{def:physics:phyx-mini-0520:phyxminiproblems-problemphyxmini0520-wavenumberininversemeters}
  \lean{PhyXMiniProblems.ProblemPhyXMini0520.waveNumberInInverseMeters}
  \uses{def:physics:phyx-mini-0520:phyxminiproblems-problemphyxmini0520-wavenumberquantity}
  Inverse-metre readout of a physical wave number
\end{definition}
\begin{proof}
  This is the declared model definition of wave Number In Inverse Meters.
\end{proof}
\begin{definition}[Particle Species]
  \label{def:physics:phyx-mini-0520:phyxminiproblems-problemphyxmini0520-particlespecies}
  \lean{PhyXMiniProblems.ProblemPhyXMini0520.ParticleSpecies}
  Particle species used in the simple nuclear model
\end{definition}
\begin{proof}
  This is the declared model definition of Particle Species.
\end{proof}
\begin{definition}[Nuclear Potential Model]
  \label{def:physics:phyx-mini-0520:phyxminiproblems-problemphyxmini0520-nuclearpotentialmodel}
  \lean{PhyXMiniProblems.ProblemPhyXMini0520.NuclearPotentialModel}
  Qualitative model selected for the nuclear potential
\end{definition}
\begin{proof}
  This is the declared model definition of Nuclear Potential Model.
\end{proof}
\begin{definition}[Plot Axis]
  \label{def:physics:phyx-mini-0520:phyxminiproblems-problemphyxmini0520-plotaxis}
  \lean{PhyXMiniProblems.ProblemPhyXMini0520.PlotAxis}
  The two axes drawn in the supplied barrier graph
\end{definition}
\begin{proof}
  This is the declared model definition of Plot Axis.
\end{proof}
\begin{definition}[Plot Axis Quantity]
  \label{def:physics:phyx-mini-0520:phyxminiproblems-problemphyxmini0520-plotaxisquantity}
  \lean{PhyXMiniProblems.ProblemPhyXMini0520.PlotAxisQuantity}
  Physical quantity assigned to an axis of the supplied graph
\end{definition}
\begin{proof}
  This is the declared model definition of Plot Axis Quantity.
\end{proof}
\begin{definition}[Potential Curve Shape]
  \label{def:physics:phyx-mini-0520:phyxminiproblems-problemphyxmini0520-potentialcurveshape}
  \lean{PhyXMiniProblems.ProblemPhyXMini0520.PotentialCurveShape}
  Shape of the potential-energy curve in the supplied graph
\end{definition}
\begin{proof}
  This is the declared model definition of Potential Curve Shape.
\end{proof}
\begin{definition}[Figure Annotation]
  \label{def:physics:phyx-mini-0520:phyxminiproblems-problemphyxmini0520-figureannotation}
  \lean{PhyXMiniProblems.ProblemPhyXMini0520.FigureAnnotation}
  Literal symbolic or numerical annotations visible in the primary image
\end{definition}
\begin{proof}
  This is the declared model definition of Figure Annotation.
\end{proof}
\begin{definition}[Rectangular Barrier Figure]
  \label{def:physics:phyx-mini-0520:phyxminiproblems-problemphyxmini0520-rectangularbarrierfigure}
  \lean{PhyXMiniProblems.ProblemPhyXMini0520.RectangularBarrierFigure}
  \uses{def:physics:phyx-mini-0520:phyxminiproblems-problemphyxmini0520-lengthquantity, def:physics:phyx-mini-0520:phyxminiproblems-problemphyxmini0520-energyquantity, def:physics:phyx-mini-0520:phyxminiproblems-problemphyxmini0520-plotaxis, def:physics:phyx-mini-0520:phyxminiproblems-problemphyxmini0520-plotaxisquantity, def:physics:phyx-mini-0520:phyxminiproblems-problemphyxmini0520-potentialcurveshape, def:physics:phyx-mini-0520:phyxminiproblems-problemphyxmini0520-figureannotation}
  The physical width and energy levels represented by the graph are fields, separate from the Boolean record of which labels are printed
\end{definition}
\begin{proof}
  This is the declared model definition of Rectangular Barrier Figure.
\end{proof}
\begin{definition}[Alpha Rectangular Barrier Setup]
  \label{def:physics:phyx-mini-0520:phyxminiproblems-problemphyxmini0520-alpharectangularbarriersetup}
  \lean{PhyXMiniProblems.ProblemPhyXMini0520.AlphaRectangularBarrierSetup}
  \uses{def:physics:phyx-mini-0520:phyxminiproblems-problemphyxmini0520-lengthquantity, def:physics:phyx-mini-0520:phyxminiproblems-problemphyxmini0520-massquantity, def:physics:phyx-mini-0520:phyxminiproblems-problemphyxmini0520-energyquantity, def:physics:phyx-mini-0520:phyxminiproblems-problemphyxmini0520-actionquantity, def:physics:phyx-mini-0520:phyxminiproblems-problemphyxmini0520-wavenumberquantity, def:physics:phyx-mini-0520:phyxminiproblems-problemphyxmini0520-particlespecies, def:physics:phyx-mini-0520:phyxminiproblems-problemphyxmini0520-nuclearpotentialmodel, def:physics:phyx-mini-0520:phyxminiproblems-problemphyxmini0520-rectangularbarrierfigure}
  Independent physical quantities in the alpha-tunneling model. In particular, the attenuation wave number, exact tunneling probability, and opaque-barrier estimate are independent fields. The governing laws below constrain them without defining any one of them to be a recorded answer value
\end{definition}
\begin{proof}
  This is the declared model definition of Alpha Rectangular Barrier Setup.
\end{proof}
\begin{definition}[Matches Alpha Nucleus Scenario]
  \label{def:physics:phyx-mini-0520:phyxminiproblems-problemphyxmini0520-matchesalphanucleusscenario}
  \lean{PhyXMiniProblems.ProblemPhyXMini0520.MatchesAlphaNucleusScenario}
  \uses{def:physics:phyx-mini-0520:phyxminiproblems-problemphyxmini0520-alpharectangularbarriersetup}
  Qualitative physical scenario stated in the problem prose
\end{definition}
\begin{proof}
  This is the declared model definition of Matches Alpha Nucleus Scenario.
\end{proof}
\begin{definition}[Matches Supplied Barrier Figure]
  \label{def:physics:phyx-mini-0520:phyxminiproblems-problemphyxmini0520-matchessuppliedbarrierfigure}
  \lean{PhyXMiniProblems.ProblemPhyXMini0520.MatchesSuppliedBarrierFigure}
  \uses{def:physics:phyx-mini-0520:phyxminiproblems-problemphyxmini0520-alpharectangularbarriersetup}
  Axes, labels, rectangular geometry, and energy levels read from the primary image. This premise contains no transmission probability or answer choice
\end{definition}
\begin{proof}
  This is the declared model definition of Matches Supplied Barrier Figure.
\end{proof}
\begin{definition}[Matches Problem Readouts]
  \label{def:physics:phyx-mini-0520:phyxminiproblems-problemphyxmini0520-matchesproblemreadouts}
  \lean{PhyXMiniProblems.ProblemPhyXMini0520.MatchesProblemReadouts}
  \uses{def:physics:phyx-mini-0520:phyxminiproblems-problemphyxmini0520-lengthinfemtometers, def:physics:phyx-mini-0520:phyxminiproblems-problemphyxmini0520-massinkilograms, def:physics:phyx-mini-0520:phyxminiproblems-problemphyxmini0520-energyinmegaelectronvolts, def:physics:phyx-mini-0520:phyxminiproblems-problemphyxmini0520-alpharectangularbarriersetup}
  Numerical physical data in the prose and figure. The final field states that the incident level is the displayed deficit below the barrier top; it does not state a transmission probability
\end{definition}
\begin{proof}
  This is the declared model definition of Matches Problem Readouts.
\end{proof}
\begin{definition}[Uses Standard Reduced Planck Constant]
  \label{def:physics:phyx-mini-0520:phyxminiproblems-problemphyxmini0520-usesstandardreducedplanckconstant}
  \lean{PhyXMiniProblems.ProblemPhyXMini0520.UsesStandardReducedPlanckConstant}
  \uses{def:physics:phyx-mini-0520:phyxminiproblems-problemphyxmini0520-actioninjouleseconds, def:physics:phyx-mini-0520:phyxminiproblems-problemphyxmini0520-alpharectangularbarriersetup}
  The standard reduced Planck constant used by the numerical model
\end{definition}
\begin{proof}
  This is the declared model definition of Uses Standard Reduced Planck Constant.
\end{proof}
\begin{definition}[Matches Physlib Rectangular Barrier]
  \label{def:physics:phyx-mini-0520:phyxminiproblems-problemphyxmini0520-matchesphyslibrectangularbarrier}
  \lean{PhyXMiniProblems.ProblemPhyXMini0520.MatchesPhyslibRectangularBarrier}
  \uses{def:physics:phyx-mini-0520:phyxminiproblems-problemphyxmini0520-lengthinmeters, def:physics:phyx-mini-0520:phyxminiproblems-problemphyxmini0520-massinkilograms, def:physics:phyx-mini-0520:phyxminiproblems-problemphyxmini0520-energyinjoules, def:physics:phyx-mini-0520:phyxminiproblems-problemphyxmini0520-alpharectangularbarriersetup}
  Calibrate Physlib's scalar rectangular-barrier parameters by coherent-SI readouts of the dimensionful physical quantities. The absolute radial location of the two faces is left free because the image specifies only their separation
\end{definition}
\begin{proof}
  This is the declared model definition of Matches Physlib Rectangular Barrier.
\end{proof}
\begin{definition}[Has Physical Tunneling Parameters]
  \label{def:physics:phyx-mini-0520:phyxminiproblems-problemphyxmini0520-hasphysicaltunnelingparameters}
  \lean{PhyXMiniProblems.ProblemPhyXMini0520.HasPhysicalTunnelingParameters}
  \uses{def:physics:phyx-mini-0520:phyxminiproblems-problemphyxmini0520-lengthinmeters, def:physics:phyx-mini-0520:phyxminiproblems-problemphyxmini0520-massinkilograms, def:physics:phyx-mini-0520:phyxminiproblems-problemphyxmini0520-energyinjoules, def:physics:phyx-mini-0520:phyxminiproblems-problemphyxmini0520-actioninjouleseconds, def:physics:phyx-mini-0520:phyxminiproblems-problemphyxmini0520-wavenumberininversemeters, def:physics:phyx-mini-0520:phyxminiproblems-problemphyxmini0520-alpharectangularbarriersetup}
  Positivity and range conditions selecting a physical tunneling setup
\end{definition}
\begin{proof}
  This is the declared model definition of Has Physical Tunneling Parameters.
\end{proof}
\begin{definition}[barrier Opacity]
  \label{def:physics:phyx-mini-0520:phyxminiproblems-problemphyxmini0520-barrieropacity}
  \lean{PhyXMiniProblems.ProblemPhyXMini0520.barrierOpacity}
  \uses{def:physics:phyx-mini-0520:phyxminiproblems-problemphyxmini0520-lengthinmeters, def:physics:phyx-mini-0520:phyxminiproblems-problemphyxmini0520-wavenumberininversemeters, def:physics:phyx-mini-0520:phyxminiproblems-problemphyxmini0520-alpharectangularbarriersetup}
  The dimensionless barrier opacity κ L at the supplied finite width
\end{definition}
\begin{proof}
  This is the declared model definition of barrier Opacity.
\end{proof}
\begin{definition}[exact Finite Barrier Transmission]
  \label{def:physics:phyx-mini-0520:phyxminiproblems-problemphyxmini0520-exactfinitebarriertransmission}
  \lean{PhyXMiniProblems.ProblemPhyXMini0520.exactFiniteBarrierTransmission}
  The exact transmission coefficient for a one-dimensional rectangular barrier with equal exterior potentials, written as a function of the dimensionless opacity x = κ L. The physical branch assumptions 0 < E < U₀ are imposed where this function is used as a law
\end{definition}
\begin{proof}
  This is the declared model definition of exact Finite Barrier Transmission.
\end{proof}
\begin{definition}[opaque Barrier Transmission Estimate]
  \label{def:physics:phyx-mini-0520:phyxminiproblems-problemphyxmini0520-opaquebarriertransmissionestimate}
  \lean{PhyXMiniProblems.ProblemPhyXMini0520.opaqueBarrierTransmissionEstimate}
  The leading opaque-barrier estimate, including the interface-matching prefactor. This is a named estimator, not the exact finite-opacity probability
\end{definition}
\begin{proof}
  This is the declared model definition of opaque Barrier Transmission Estimate.
\end{proof}
\begin{lemma}[exact Finite Barrier Transmission is Equivalent opaque Estimate]
  \label{lem:physics:phyx-mini-0520:phyxminiproblems-problemphyxmini0520-exactfinitebarriertransmission-isequivalent-opaqueestimate}
  \lean{PhyXMiniProblems.ProblemPhyXMini0520.exactFiniteBarrierTransmission_isEquivalent_opaqueEstimate}
  \uses{def:physics:phyx-mini-0520:phyxminiproblems-problemphyxmini0520-exactfinitebarriertransmission, def:physics:phyx-mini-0520:phyxminiproblems-problemphyxmini0520-opaquebarriertransmissionestimate}
  For fixed energies 0 < E < U₀, the exact rectangular-barrier coefficient is asymptotically equivalent to the opaque-barrier estimate as κ L → +∞. This is the precise large-opacity contract behind the approximation
\end{lemma}
\begin{proof}
  The conclusion follows from the governing relations and figure readouts named in the statement of exact Finite Barrier Transmission is Equivalent opaque Estimate.
\end{proof}
\begin{definition}[Obeys Finite Rectangular Barrier Tunneling]
  \label{def:physics:phyx-mini-0520:phyxminiproblems-problemphyxmini0520-obeysfiniterectangularbarriertunneling}
  \lean{PhyXMiniProblems.ProblemPhyXMini0520.ObeysFiniteRectangularBarrierTunneling}
  \uses{def:physics:phyx-mini-0520:phyxminiproblems-problemphyxmini0520-massinkilograms, def:physics:phyx-mini-0520:phyxminiproblems-problemphyxmini0520-energyinjoules, def:physics:phyx-mini-0520:phyxminiproblems-problemphyxmini0520-actioninjouleseconds, def:physics:phyx-mini-0520:phyxminiproblems-problemphyxmini0520-wavenumberininversemeters, def:physics:phyx-mini-0520:phyxminiproblems-problemphyxmini0520-alpharectangularbarriersetup, def:physics:phyx-mini-0520:phyxminiproblems-problemphyxmini0520-barrieropacity, def:physics:phyx-mini-0520:phyxminiproblems-problemphyxmini0520-exactfinitebarriertransmission, def:physics:phyx-mini-0520:phyxminiproblems-problemphyxmini0520-opaquebarriertransmissionestimate}
  The exact finite-barrier physics and the separately named opaque-limit estimate used for the recorded multiple-choice computation. In the classically forbidden region, ℏ² κ² = 2 m (U₀ - E). The physical finite-width probability is the exact hyperbolic-sine coefficient. The opaque estimate is the leading asymptotic expression 16 (E/U₀) (1-E/U₀) exp(-2 κ L). Every term is an SI readout of a dimensionally typed quantity except the dimensionless ratios, exponent, and probability.
\end{definition}
\begin{proof}
  This is the declared model definition of Obeys Finite Rectangular Barrier Tunneling.
\end{proof}
\begin{definition}[Answer Choice]
  \label{def:physics:phyx-mini-0520:phyxminiproblems-problemphyxmini0520-answerchoice}
  \lean{PhyXMiniProblems.ProblemPhyXMini0520.AnswerChoice}
  Labels of the four tunneling-probability choices
\end{definition}
\begin{proof}
  This is the declared model definition of Answer Choice.
\end{proof}
\begin{definition}[displayed Tunneling Probability]
  \label{def:physics:phyx-mini-0520:phyxminiproblems-problemphyxmini0520-displayedtunnelingprobability}
  \lean{PhyXMiniProblems.ProblemPhyXMini0520.displayedTunnelingProbability}
  \uses{def:physics:phyx-mini-0520:phyxminiproblems-problemphyxmini0520-answerchoice}
  Dimensionless probabilities printed beside the four choices
\end{definition}
\begin{proof}
  This is the declared model definition of displayed Tunneling Probability.
\end{proof}
\begin{definition}[recorded Dataset Answer]
  \label{def:physics:phyx-mini-0520:phyxminiproblems-problemphyxmini0520-recordeddatasetanswer}
  \lean{PhyXMiniProblems.ProblemPhyXMini0520.recordedDatasetAnswer}
  \uses{def:physics:phyx-mini-0520:phyxminiproblems-problemphyxmini0520-answerchoice}
  The answer label recorded by the source dataset
\end{definition}
\begin{proof}
  This is the declared model definition of recorded Dataset Answer.
\end{proof}
\begin{definition}[Rounds To Three Decimal Places]
  \label{def:physics:phyx-mini-0520:phyxminiproblems-problemphyxmini0520-roundstothreedecimalplaces}
  \lean{PhyXMiniProblems.ProblemPhyXMini0520.RoundsToThreeDecimalPlaces}
  value rounds to reported at three decimal places. The half-open interval implements ordinary nearest-thousandth rounding away from the upper tie
\end{definition}
\begin{proof}
  This is the declared model definition of Rounds To Three Decimal Places.
\end{proof}
\begin{definition}[Matches Displayed Tunneling Probability]
  \label{def:physics:phyx-mini-0520:phyxminiproblems-problemphyxmini0520-matchesdisplayedtunnelingprobability}
  \lean{PhyXMiniProblems.ProblemPhyXMini0520.MatchesDisplayedTunnelingProbability}
  \uses{def:physics:phyx-mini-0520:phyxminiproblems-problemphyxmini0520-answerchoice, def:physics:phyx-mini-0520:phyxminiproblems-problemphyxmini0520-displayedtunnelingprobability, def:physics:phyx-mini-0520:phyxminiproblems-problemphyxmini0520-roundstothreedecimalplaces}
  A displayed choice agrees with a supplied probability to three decimals
\end{definition}
\begin{proof}
  This is the declared model definition of Matches Displayed Tunneling Probability.
\end{proof}
\begin{definition}[Is Unique Matching Tunneling Choice]
  \label{def:physics:phyx-mini-0520:phyxminiproblems-problemphyxmini0520-isuniquematchingtunnelingchoice}
  \lean{PhyXMiniProblems.ProblemPhyXMini0520.IsUniqueMatchingTunnelingChoice}
  \uses{def:physics:phyx-mini-0520:phyxminiproblems-problemphyxmini0520-answerchoice, def:physics:phyx-mini-0520:phyxminiproblems-problemphyxmini0520-matchesdisplayedtunnelingprobability}
  A choice is the unique displayed value matching the modeled probability
\end{definition}
\begin{proof}
  This is the declared model definition of Is Unique Matching Tunneling Choice.
\end{proof}
\begin{definition}[No Displayed Tunneling Choice Matches]
  \label{def:physics:phyx-mini-0520:phyxminiproblems-problemphyxmini0520-nodisplayedtunnelingchoicematches}
  \lean{PhyXMiniProblems.ProblemPhyXMini0520.NoDisplayedTunnelingChoiceMatches}
  \uses{def:physics:phyx-mini-0520:phyxminiproblems-problemphyxmini0520-answerchoice, def:physics:phyx-mini-0520:phyxminiproblems-problemphyxmini0520-matchesdisplayedtunnelingprobability}
  No source choice rounds to the supplied probability at three decimals
\end{definition}
\begin{proof}
  This is the declared model definition of No Displayed Tunneling Choice Matches.
\end{proof}
\begin{lemma}[incident kinetic energy megaelectronvolts eq]
  \label{lem:physics:phyx-mini-0520:phyxminiproblems-problemphyxmini0520-incident-kinetic-energy-megaelectronvolts-eq}
  \lean{PhyXMiniProblems.ProblemPhyXMini0520.incident_kinetic_energy_megaelectronvolts_eq}
  \uses{def:physics:phyx-mini-0520:phyxminiproblems-problemphyxmini0520-energyinmegaelectronvolts, def:physics:phyx-mini-0520:phyxminiproblems-problemphyxmini0520-alpharectangularbarriersetup, def:physics:phyx-mini-0520:phyxminiproblems-problemphyxmini0520-matchesproblemreadouts}
  The displayed energy separation implies an incident energy of 29 MeV
\end{lemma}
\begin{proof}
  The conclusion follows from the governing relations and figure readouts named in the statement of incident kinetic energy megaelectronvolts eq.
\end{proof}
\begin{lemma}[barrier opacity rounds to 0 875]
  \label{lem:physics:phyx-mini-0520:phyxminiproblems-problemphyxmini0520-barrier-opacity-rounds-to-0-875}
  \lean{PhyXMiniProblems.ProblemPhyXMini0520.barrier_opacity_rounds_to_0_875}
  \uses{def:physics:phyx-mini-0520:phyxminiproblems-problemphyxmini0520-alpharectangularbarriersetup, def:physics:phyx-mini-0520:phyxminiproblems-problemphyxmini0520-matchesproblemreadouts, def:physics:phyx-mini-0520:phyxminiproblems-problemphyxmini0520-usesstandardreducedplanckconstant, def:physics:phyx-mini-0520:phyxminiproblems-problemphyxmini0520-hasphysicaltunnelingparameters, def:physics:phyx-mini-0520:phyxminiproblems-problemphyxmini0520-barrieropacity, def:physics:phyx-mini-0520:phyxminiproblems-problemphyxmini0520-obeysfiniterectangularbarriertunneling, def:physics:phyx-mini-0520:phyxminiproblems-problemphyxmini0520-roundstothreedecimalplaces}
  The source data put the opacity at about 0.875, so this finite barrier is not in a quantitatively large-opacity regime
\end{lemma}
\begin{proof}
  The conclusion follows from the governing relations and figure readouts named in the statement of barrier opacity rounds to 0 875.
\end{proof}
\begin{lemma}[exact tunneling probability rounds to 0 116]
  \label{lem:physics:phyx-mini-0520:phyxminiproblems-problemphyxmini0520-exact-tunneling-probability-rounds-to-0-116}
  \lean{PhyXMiniProblems.ProblemPhyXMini0520.exact_tunneling_probability_rounds_to_0_116}
  \uses{def:physics:phyx-mini-0520:phyxminiproblems-problemphyxmini0520-alpharectangularbarriersetup, def:physics:phyx-mini-0520:phyxminiproblems-problemphyxmini0520-matchesproblemreadouts, def:physics:phyx-mini-0520:phyxminiproblems-problemphyxmini0520-usesstandardreducedplanckconstant, def:physics:phyx-mini-0520:phyxminiproblems-problemphyxmini0520-hasphysicaltunnelingparameters, def:physics:phyx-mini-0520:phyxminiproblems-problemphyxmini0520-obeysfiniterectangularbarriertunneling, def:physics:phyx-mini-0520:phyxminiproblems-problemphyxmini0520-roundstothreedecimalplaces}
  The exact finite rectangular-barrier coefficient is about 0.116072; it rounds to 0.116, which is absent from the four displayed choices
\end{lemma}
\begin{proof}
  The conclusion follows from the governing relations and figure readouts named in the statement of exact tunneling probability rounds to 0 116.
\end{proof}
\begin{lemma}[opaque barrier estimate rounds to 0 090]
  \label{lem:physics:phyx-mini-0520:phyxminiproblems-problemphyxmini0520-opaque-barrier-estimate-rounds-to-0-090}
  \lean{PhyXMiniProblems.ProblemPhyXMini0520.opaque_barrier_estimate_rounds_to_0_090}
  \uses{def:physics:phyx-mini-0520:phyxminiproblems-problemphyxmini0520-alpharectangularbarriersetup, def:physics:phyx-mini-0520:phyxminiproblems-problemphyxmini0520-matchesproblemreadouts, def:physics:phyx-mini-0520:phyxminiproblems-problemphyxmini0520-usesstandardreducedplanckconstant, def:physics:phyx-mini-0520:phyxminiproblems-problemphyxmini0520-hasphysicaltunnelingparameters, def:physics:phyx-mini-0520:phyxminiproblems-problemphyxmini0520-obeysfiniterectangularbarriertunneling, def:physics:phyx-mini-0520:phyxminiproblems-problemphyxmini0520-roundstothreedecimalplaces}
  The leading opaque-barrier estimate is about 0.089626; unlike the exact finite-barrier probability, this estimate rounds to the source value 0.090
\end{lemma}
\begin{proof}
  The conclusion follows from the governing relations and figure readouts named in the statement of opaque barrier estimate rounds to 0 090.
\end{proof}
% --- Archon declaration topology end ---
% --- Archon physics formalization source end ---
